\chapter{SYSTEM DESIGN AND METHODOLOGY}

This chapter describes how the prediction system is designed and how the data flows through it. It begins with the overall architecture, then details the real-time data collection subsystem, the extraction of stop-to-stop travel-time segments, the feature engineering pipeline, and finally the architectures of the models that are trained and compared.

\section{System Architecture}

The system is organized as a five-stage pipeline in which the output of each stage is a well-defined artifact that becomes the input of the next stage. This modular design makes the pipeline reproducible and allows each stage to be tested independently.

\begin{enumerate}
    \item \textbf{Data collection:} real-time bus positions, weather, and traffic are polled from external services and stored in a relational database.
    \item \textbf{Trip and segment extraction:} raw stop-arrival events are grouped into trips and converted into stop-to-stop travel-time segments.
    \item \textbf{Feature engineering:} each segment is enriched with temporal, spatial, schedule, historical, and dwell-time features.
    \item \textbf{Model training:} baseline, classical, and deep learning models are trained on a chronological training split.
    \item \textbf{Evaluation:} all models are compared on a held-out test split using error metrics, statistical tests, and a comparison with the reference study.
\end{enumerate}

A strict separation between data collection (online) and processing (offline) is maintained: the collector only writes raw observations, while all feature construction and modeling are performed offline from the stored data. This guarantees that results can be regenerated exactly from the database at any time.

\section{Data Collection}

The data collector continuously gathers three complementary data sources for three İzmir ESHOT routes (502, 268, and 565) and persists them in an SQLite database. Three routes with different traffic profiles are collected so that the generalization of the method can later be tested.

\subsection{Real-Time Bus Positions}

The primary source is the İzmir open transit API, which returns the live coordinates of all active buses on a route. The collector polls this endpoint every 30~seconds for all three routes in parallel. For each returned bus, the nearest stop on the route is determined using the Haversine great-circle distance between the bus and every stop of the relevant direction:
\begin{equation}
d = 2R \arcsin\!\sqrt{\sin^{2}\!\Big(\tfrac{\Delta\phi}{2}\Big) + \cos\phi_{1}\cos\phi_{2}\,\sin^{2}\!\Big(\tfrac{\Delta\lambda}{2}\Big)},
\end{equation}
where $\phi$ and $\lambda$ are latitude and longitude in radians and $R$ is the Earth's radius. When a bus first comes within 150~m of a stop, a stop-arrival event is recorded. The stop coordinates and sequences themselves are loaded from the static GTFS feed, so that all routes share a single authoritative source of stop information.

\subsection{Weather and Traffic}

Weather is collected hourly from the OpenWeatherMap service (temperature, humidity, precipitation, wind speed, visibility, and a categorical condition); when no API key is available, a deterministic mock is used so that the pipeline remains runnable. Traffic flow is sampled every 20~minutes from the TomTom service for the segments of route 502 only, because the free service tier imposes a daily request limit. Each traffic reading yields a congestion ratio (current speed divided by free-flow speed).

\subsection{Storage and Volume}

All observations are stored in an SQLite database across five tables: bus positions, stop arrivals, stop-transition (trip) events, weather readings, and traffic readings, with indexes to support efficient offline querying. Over a collection window spanning April to June 2026 (73 days with recorded data), approximately four million bus-position records were gathered across the three routes, which were later converted into tens of thousands of clean travel-time segments per route (for example, more than $81{,}000$ segments for route 502).

\section{Trip and Segment Extraction}

The stored stop-arrival events are converted into the supervised learning target offline. For each bus and travel direction, the chronologically ordered events are grouped into individual trips: a new trip is started whenever there is a time gap larger than 30~minutes or the stop sequence resets, both of which indicate the end of one service run and the beginning of another. Within a trip, the travel time of a segment is the time difference between the arrivals at two consecutive stops.

To remove noise and physically implausible values, only segments whose travel time falls between $0.33$ and $15$~minutes are retained; values outside this range correspond to GPS glitches, long layovers, or matching errors. The result is a clean table of segments, each identified by its origin and destination stop, its trip, its timestamp, and its observed travel time, which is the prediction target.

\section{Feature Engineering}

Each segment is described by a set of features grouped into five families, summarized in Table~\ref{tab:features}. The features were developed in successive iterations and then reduced to a compact, high-signal set through the feature-selection experiments described in Chapter~4.

\begin{table}[!htbp]
    \centering
    \caption{\label{tab:features}Feature groups used for inter-stop travel-time prediction.}
    \begin{tabular}{@{}lp{8.2cm}@{}}
        \toprule
        \textbf{Group} & \textbf{Features} \\
        \midrule
        Temporal & hour, day of week, weekend flag (cyclic sine/cosine encodings for the LSTM) \\
        Spatial & inter-stop distance, progress along the route \\
        Schedule (GTFS) & scheduled travel time, deviation from schedule \\
        Lag / deviation & previous segment travel time and deviation, cumulative and rolling-window deviation within the trip \\
        Historical & per-stop historical median travel time and ratio, previous segment speed \\
        Dwell & dwell time at the stop and the previous dwell time \\
        Cold-start & trip-start indicator \\
        \bottomrule
    \end{tabular}
\end{table}

\subsection{Schedule and Dwell Features}

Two feature families are central to this project's hypotheses. The \textbf{schedule} features are derived from the GTFS \texttt{stop\_times} file: the scheduled travel time provides a strong baseline expectation for each segment, and the deviation of the observed time from it captures how the current trip is running relative to plan. The \textbf{dwell-time} features are computed from the raw GPS records: the time a bus spends within 50~m of a stop (bounded between 10 and 600~seconds) approximates the boarding and alighting duration, which strongly influences the following segment.

\subsection{Leakage-Safe Construction}

Because the target is a travel time and several features are derived from past travel times, care is taken to avoid information leakage. All lag features use only values from \emph{earlier} segments of the same trip, and all historical statistics (such as the per-stop median) are computed exclusively from the training split. The data are split chronologically into $80\%$ training and $20\%$ testing, so that the model is always evaluated on the temporally most recent data and never sees future information during training.

\subsection{Target Transformation}

The travel-time distribution is right-skewed and strictly positive. To stabilize it, the target is modeled in logarithmic space using the transformation
\begin{equation}
y' = \ln(1 + y), \qquad \hat{y} = e^{\hat{y}'} - 1,
\end{equation}
where $y$ is the travel time in minutes. Predictions are transformed back to the original scale and clipped to non-negative values. This transformation reduces the influence of large outliers and improves the percentage-error behavior of the models.

\section{Model Architectures}

Six predictors of increasing complexity are trained and compared.

\subsection{Baseline Models}

Two non-learned baselines establish a lower bound on acceptable performance. The \textbf{Naïve (GTFS)} predictor simply returns the scheduled travel time. The \textbf{Historical Average} predictor returns the mean observed travel time for the matching combination of hour, weekend flag, and origin stop, computed on the training split. A \textbf{Linear Regression} model provides a simple learned baseline.

\subsection{Classical Machine Learning Models}

The \textbf{Random Forest} uses $300$ trees with a maximum depth of $15$, and the \textbf{XGBoost} model uses $500$ trees with a maximum depth of $8$, a learning rate of $0.03$, and subsampling for regularization. Both are trained on the log-transformed target and on the compact feature set. These two models represent the strongest classical predictors in the study.

\subsection{Dual-Input LSTM}

The deep learning model is a dual-input network designed to combine sequential and static information, illustrated conceptually in Figure~\ref{fig:lstm}. The \textbf{sequence branch} receives a sliding window of the most recent segments of the current trip, each represented by its lagged travel time, scheduled time, distance, and route progress; this window is processed by a two-layer LSTM with $128$ hidden units and dropout. The \textbf{context branch} receives the static features of the segment to be predicted (temporal, weather, historical, dwell, and the trip-start indicator) through a fully connected layer with $32$ units and a ReLU activation. The two branch outputs are concatenated and passed through a $64$-unit hidden layer to a single linear output.

The network is trained on the log-transformed target with the Huber loss, which behaves quadratically for small errors and linearly for large ones, making the model robust to occasional large outliers. Optimization uses Adam with early stopping on a validation split. The size of the input window is treated as a hyperparameter and selected experimentally in Chapter~4.

\begin{figure}[!htbp]
    \centering
    \includegraphics[width=0.7\textwidth]{frog.jpg}
    \caption{\label{fig:lstm}Conceptual structure of the dual-input LSTM: a sequence branch over recent segments and a context branch over static features are concatenated to produce the travel-time prediction. (Placeholder image to be replaced with the final architecture diagram.)}
\end{figure}
